%%%%%%%%%%%%%%%%%%%%%%%%%%%%%%%%%%%%%%%%%%%%%%%%%
\section{Dataset Sourcing and Class Taxonomy}\label{sec:data_sourcing_and_taxonomy}

Before any imagery could be gathered, the set of species the detector was expected to recognize had to be fixed, since it determines both what counts as a usable image and how large a corpus is realistically obtainable. This section describes how the class taxonomy was derived from the \textit{SpeciesNet} taxonomy \cite{gadotCropNotCrop2024a}, how candidate image sources were assessed for licensing compatibility and suitability, and how the sourcing effort was executed. This included a substantial detour through the \textit{LILA BC} camera-trap archives, which was ultimately abandoned once its imagery proved a poor match for the deployment scenario.

%%%%%%%%%%%%%%%%%%%%%%%%%%%%%%%%%%%%%%%%%%%%%%%%%
\subsection{Target Class Definition}\label{sec:target_class_definition}

The class universe originates from the taxonomy used by \textit{SpeciesNet} \cite{gadotCropNotCrop2024a}, a species-classification model published by Google and trained in cooperation with the wildlife-camera-trap community. The full taxonomy comprises $3{,}537$ labels, of which approximately $1{,}415$ correspond to mammal taxa at varying levels of specificity (species, genus, family, and higher clades). Since the AX Visio binocular is intended for non-bird mammal observation, the taxonomy was first restricted to this subset. The AX Visio product owner then manually filtered this subset down to $483$ candidate labels using a dedicated filtering script, removing taxa that are not realistically observable in the product's target markets or that are duplicated across taxonomic levels.  % resources/refineLabels.py

The criterion applied at this stage is worth making explicit, because it is an observability criterion rather than a taxonomic one: the filter reduces the mammal taxonomy to \emph{diurnal land mammals that are likely to be photographed}. Its exclusions follow directly from that premise. Nocturnal taxa (bats, bush babies, several marsupial groups) are removed because the product is used in daylight. Range-limited endemics are removed because a user in the target markets will not encounter them. Small rodents are removed wholesale as implausible binocular subjects. Subspecies are collapsed into their parent species. Visually inseparable congeners are merged rather than dropped, on the grounds that the group remains recognizable even where its members are not. A second criterion governs additions: species that the automatic filter had removed but that a reviewer judged commonly encountered were added back individually. The final list is therefore the product of an expert pass in both directions, rather than of a rule applied blindly.

A separate, quantitative criterion governed whether a species could realistically be supported with imagery at all. Per-taxon image counts were surveyed on \textit{GBIF} \cite{GBIF} and logged in a dedicated image-count table, and a count of at least $300$ images was used as the inclusion threshold.  % data/gbif/metadata/GBIF_image_counts_v1.csv
Because the survey caps at approximately $300$ observations per taxon ($125$ of $479$ surveyed taxa sit exactly at the ceiling), a count at the cap means only \enquote{at least $300$}. A low count, by contrast, is evidence of genuine scarcity rather than a sampling artifact. This asymmetry is what makes the measure usable as a rarity proxy later, in \Cref{sec:generator_comparison_classes}.

A class list of this size is, however, not compatible with the computational budget of the target hardware. A review of nano-scale detection architectures suitable for the \textit{Qualcomm QCS605} found consistent evidence that classification and detection accuracy degrades markedly once the number of output classes exceeds roughly $200$--$300$. The architectures reviewed include \textit{YOLOv8n} at $3.2\,\mathrm{M}$ parameters, \textit{YOLOv10n} at $2.7\,\mathrm{M}$, \textit{YOLOv11n} at $2.6\,\mathrm{M}$, and \textit{MobileNetV3-Small} at $1.5\,\mathrm{M}$. The degradation follows because the discriminative capacity of a $1$--$3\,\mathrm{M}$-parameter backbone is spread too thinly across a large, fine-grained label space. This established an approximate ceiling of $200$--$250$ classes for the primary experiment. This ceiling should be read as a design heuristic adopted for this project, rather than as an established empirical result. It was derived from a survey of applied edge-deployment literature and vendor documentation, rather than from a controlled study. No sweep over class count was run within this thesis to confirm it. It is stated here as a premise of the dataset design, in the same spirit as the student-architecture choice discussed in \Cref{sec:model_selection}.

Three scope exclusions apply on top of the class budget, each for a reason specific to the deployment scenario. Birds are excluded because they are served by a separate model in the product. Marine mammals are excluded except for pinnipeds, which are observable from shore. Bats are excluded as nocturnal. Domestic animals, by contrast, are deliberately included despite consuming class budget that could have gone to wild species. A detector that misidentifies a domestic dog as a coyote undermines user confidence more than one that fails to recognize an uncommon species at all.

The final class list settled on $225$ classes, structured across three taxonomic tiers. This structure keeps the fine-grained species count within budget while still covering visually and ecologically distinct groups that could not be meaningfully merged. The tiers comprise $164$ classes at species level (Tier 1), $41$ at genus level (Tier 2), and $20$ at family or higher clade level (Tier 3). Tier 2 is used where individual species within a genus are not reliably visually distinguishable. Tier 3 is used for taxa that are rare, visually homogeneous as a group, or both. \Cref{tab:class-taxonomy-orders} breaks the $225$ classes down by taxonomic order.

\begin{table}[H]
\centering
\caption{Composition of the 225-class taxonomy by mammalian order.}
\label{tab:class-taxonomy-orders}
\begin{tabular}{@{}lrr@{}}
\toprule
Order & Classes & Share \\
\midrule
Carnivora & $75$ & $33\%$ \\
Cetartiodactyla & $65$ & $29\%$ \\
Primates & $24$ & $11\%$ \\
Rodentia & $24$ & $11\%$ \\
Perissodactyla & $10$ & $4\%$ \\
Marsupials and Monotremes & $9$ & $4\%$ \\
Other & $18$ & $8\%$ \\
\midrule
Total & $225$ & $100\%$ \\
\bottomrule
\end{tabular}
\end{table}

Five principles governed selection within this budget. The expert curation embodied in the $483$-label list is respected by default, so a species is removed or merged only with a stated reason. Visual distinctiveness drives granularity: a species is kept at its own level only if a convolutional network could plausibly separate it from its relatives in a daylight image at binocular range. Taxa requiring cranial, dental or genetic examination to tell apart are consolidated upward. Expected encounter frequency in the target markets outweighs taxonomic completeness. Additions are treated as cheaper than removals, since adding one commonly encountered species costs less capacity than retaining several obscure ones. And the total is held at the class-count target rather than allowed to drift.

The tier structure follows from the second principle, but its justification is as much about the user as about the model. A user observing a chipmunk is well served by the label \enquote{chipmunk} and does not need \textit{Tamias alpinus} distinguished from \textit{T. speciosus}. Collapsing such a genus therefore costs the product almost nothing, while freeing capacity for classes where the distinction is visible and wanted. A hierarchical inference scheme was also designed at this stage, in which species, genus and family predictions would be emitted simultaneously and the displayed label would degrade gracefully as confidence fell. It is documented here for completeness but was not implemented: the model evaluated in this thesis is a flat $225$-way detector, and the graceful-degradation behavior would be a natural extension rather than a property of the system as built.

Applying these principles gives the consolidation its characteristic shape, which is heavily uneven by group. Primates were reduced by roughly $63\%$ relative to the product-owner list, mustelids by $59\%$, viverrids by $75\%$, duikers by $85\%$, and marsupials by $64\%$, while squirrels were not reduced at all. The pattern is not arbitrary. The heavily reduced groups are precisely those combining many range-limited species with low inter-species visual separability, whereas North American squirrels are common in the primary market and visually distinct. Class selection therefore favored visual distinctiveness and expected encounter frequency over strict taxonomic completeness, following the principle that additions are cheaper to correct than removals. Where a common species was missing from the initial $483$-label pool (for instance several North American squirrel species and common domestic animals), it was added in preference to retaining an obscure, rarely-photographed endemic species. Retaining such a species would consume the same class budget for negligible practical value. An extended $480$-class list was also documented as a secondary artifact. It comprises the $483$ product-owner-filtered labels, minus $2$ duplicates and $18$ removed labels, plus $17$ genus- or family-level fallback labels for coverage. It was intended for a larger student architecture or a two-stage classification pipeline. This extended list was not used for the primary knowledge-distillation experiment and is not discussed further in this chapter.

%%%%%%%%%%%%%%%%%%%%%%%%%%%%%%%%%%%%%%%%%%%%%%%%%
\subsection{Candidate Data Sources and Licensing Assessment}\label{sec:candidate_sources_licensing}

With the class list fixed, candidate image sources were surveyed against two criteria: commercial usability of the license, and suitability of the imagery for a close-to-mid-range, daylight, binocular-observation deployment scenario. Because the trained weights are intended for a commercial product, the licensing policy adopted was to exclude any dataset carrying a NonCommercial (NC), NoDerivatives (ND), or ShareAlike (SA) clause. The reasoning behind the ShareAlike exclusion in particular concerns the weights rather than the images. It is legally unsettled whether the parameters of a trained network constitute a derivative work of the images used to fit them. If a court were to hold that they do, training on ShareAlike imagery could oblige the resulting model to be released under the same copyleft terms. The exclusion is therefore a response to an asymmetric risk (a low-probability outcome with a disproportionate consequence) rather than a judgment that such imagery is unusable in principle.

The converse also holds: a license that permits training still imposes obligations. The CC-BY and CDLA-Permissive material retained in the corpus requires attribution. This is satisfied by generating a manifest of contributors, source URIs and license types alongside the trained model. The manifest is intended for inclusion in the product documentation, rather than as any change to the training process itself. \Cref{tab:candidate-sources} summarizes the sources considered.

\begin{table}[H]
\centering
\caption{Candidate image sources and their licensing verdicts.}
\label{tab:candidate-sources}
\begin{tabularx}{\textwidth}{@{}>{\raggedright\arraybackslash}p{0.22\textwidth} X >{\raggedright\arraybackslash}p{0.20\textwidth}@{}}
\toprule
Source & Notes & Verdict \\
\midrule
iNaturalist Competition datasets \cite{hornINaturalistSpeciesClassification2018} & Taxonomically ideal (2021 edition: $246$ mammal species, $68{,}917$ training images) but distributed under a non-commercial academic license & Not allowed \\
iNaturalist Open Data (via \textit{GBIF} \cite{GBIF} / AWS S3) & Same underlying photographs, filtered to observations whose contributors opted into CC0 or CC-BY & Used \\
\textit{LILA BC} \cite{LILABCLabeled} (Snapshot Serengeti, Snapshot Safari, WCS Camera Traps, Caltech Camera Traps) & CDLA-Permissive license & License allowed. Rejected on suitability grounds, see \Cref{sec:sourcing_execution} \\
Open Images V7 \cite{kuznetsovaOpenImagesDataset2020} & CC-BY 2.0 images / CC-BY 4.0 boxes. \enquote{Mammal} class covers $95{,}335$ boxed images under colloquial (non-species-level) labels & Used, with taxonomy remapping \\
COCO 2017 \cite{linMicrosoftCOCOCommon2015} & $330\mathrm{k}$ images, $10$ mammal-adjacent classes including \texttt{person} & Used, with strict filtering (\Cref{sec:sourcing_execution}) \\
iWildCam, Wildlife-71, ATRW & Contain unverifiable web-scraped imagery or images contaminated with restricted iNaturalist data & Not allowed \\
TreeOfLife-200M / \textit{BioCLIP} \cite{stevensBioCLIPVisionFoundation2024} & Raw image bundle is license-mixed (partially CC-BY-NC-ND via FathomNet). Model weights themselves are MIT-licensed & Weights used as a pseudo-labeler only, not as an image source \\
\bottomrule
\end{tabularx}
\end{table}

Source selection also fixed a second axis that the curation pipeline later depends on: how far each source's own species labels can be trusted. \textit{iNaturalist}, \textit{GBIF} and Wikipedia-derived material are treated as label-trusted, on the grounds that their labels are attached by identifiable contributors and, in the first case, subject to community verification. Open Images and the smaller stock-image source are treated as unverified, because their labels are colloquial or of unknown provenance. The distinction is not rhetorical: unverified sources contribute to a class's usable pool only where an independent classifier corroborates the label, whereas trusted sources contribute on passing the quality gates alone (\Cref{sec:data_quality_and_curation}). The asymmetry runs in the other direction too. An image from a trusted source that fails classifier verification is queued for human review rather than discarded, since the source label may well be the correct one.

Two sources merit a further note. Open Images V7 labels animals colloquially (e.g. \enquote{Elephant} rather than a species-level label), so any image drawn from it required an additional taxonomy-remapping step before it could be attributed to one of the $225$ classes. \textit{BioCLIP} \cite{stevensBioCLIPVisionFoundation2024} is an open-weight vision-language embedding model trained on the TreeOfLife-200M corpus. It was judged usable purely as a zero-shot pseudo-labeling teacher over commercially-safe unlabeled imagery, such as CC0 stock photography. This is because the MIT license on its weights does not inherit the license restrictions of its training data.

%%%%%%%%%%%%%%%%%%%%%%%%%%%%%%%%%%%%%%%%%%%%%%%%%
\subsection{Sourcing Execution and the LILA BC Rejection}\label{sec:sourcing_execution}

The initial corpus was assembled from \textit{GBIF} \cite{GBIF} observation records, classified with \textit{SpeciesNet} v$4.0.2$a \cite{gadotCropNotCrop2024a}, yielding $66{,}881$ usable images. This corpus was highly uneven across classes: $31$ of the $225$ classes had zero images, and a further $55$ classes had fewer than $10$. Closing this gap required actively sourcing additional imagery for the affected classes rather than accepting the GBIF distribution as-is.

\paragraph{The LILA BC Camera-Trap Detour}
\textit{LILA BC} (Labeled Information Library of Alexandria: Biology and Conservation) \cite{LILABCLabeled} was the first candidate investigated for filling these gaps, since its CDLA-Permissive license is commercially compatible and several of its constituent datasets are large. Three sub-datasets were scanned against the $225$-class taxonomy: \textit{Snapshot Serengeti} \cite{swansonSnapshotSerengetiHighfrequency2015} ($7.2\,\mathrm{M}$ images, $39$ of $61$ categories matched), the \textit{Snapshot Safari} 2024 expansion ($4.0\,\mathrm{M}$ images, $62$ of $131$ categories matched), and the \textit{WCS Camera Traps} dataset \cite{lilawpWCSCameraTraps2019} ($1.4\,\mathrm{M}$ images, $137$ of $676$ categories matched). \Cref{tab:lila-bc-filtering} summarizes the result of this scan.

\begin{table}[H]
\centering
\caption{LILA BC sub-datasets scanned against the 225-class taxonomy.}
\label{tab:lila-bc-filtering}
\begin{tabular}{@{}lrrr@{}}
\toprule
Sub-dataset & Total images & Matched \& filtered & With bounding boxes \\
\midrule
Snapshot Serengeti & $7.2\,\mathrm{M}$ & $1{,}637{,}554$ & $62{,}540$ \\
Snapshot Safari (2024 expansion) & $4.0\,\mathrm{M}$ & $1{,}140{,}640$ & $0$ \\
WCS Camera Traps & $1.4\,\mathrm{M}$ & $435{,}909$ & $110{,}069$ \\
\midrule
Total & $12.6\,\mathrm{M}$ & $3{,}214{,}103$ & $172{,}609$ \\
\bottomrule
\end{tabular}
\end{table}

Combined, these three sub-datasets covered $115$ of the $225$ target classes, which on paper made \textit{LILA BC} the single most productive candidate source for closing the coverage gap. Visual sampling of the matched images, however, revealed that the imagery was a poor fit for the deployment scenario. Camera-trap photographs are predominantly long-range, frequently occluded by vegetation, affected by motion blur from triggered captures, and in many cases captured under infrared illumination at night. The \textit{AX Visio} binocular, by contrast, is used for close-to-mid-range daylight observation, so a detector trained predominantly on camera-trap imagery would learn a systematically different appearance distribution than the one it will encounter in deployment. This mismatch outweighed the sourcing convenience the corpus offered, and \textit{LILA BC} was consequently removed from the sourcing pipeline entirely. This is a negative result documented here for methodological transparency, since $12.6\,\mathrm{M}$ scanned images did not translate into any images retained.

\paragraph{Gap-Filling Sources Actually Used}
With \textit{LILA BC} excluded, three sources were used instead to close the remaining class-coverage gaps.

First, a metadata-level match was performed between the $225$-class taxonomy and \textit{iNaturalist} Open Data observation records exposed through \textit{GBIF}: $7{,}158$ taxa matched to $221$ of the $225$ classes, covering $6{,}816{,}962$ photographs in total. Of these, only $831{,}050$ ($12\%$) carried a commercially-safe license (CC0 or CC-BY). The remaining $82\%$ were CC-BY-NC and had to be excluded under the licensing policy established in \Cref{sec:candidate_sources_licensing}. Coverage after this filtering step remained uneven. The best-covered class, white-tailed deer, retained $61{,}865$ commercially-safe images, while $9$ classes (aardwolf, American mink, black-backed jackal, dingo, domestic goat, domestic pig, human, pinniped clade, and red-necked wallaby) had zero matching images and required separate sourcing.

Second, \texttt{human} imagery (needed as a negative/contextual class rather than a wildlife species) was deliberately excluded from the iNaturalist and GBIF exports for privacy reasons. Open Images' \enquote{Person} category was considered as an alternative, but was found to be dominated by portrait-style photography unrepresentative of an outdoor observation scenario. \textit{COCO} 2017 \cite{linMicrosoftCOCOCommon2015} was used instead, subject to a strict filter developed specifically for this class. Images with vehicle or animal supercategory annotations were excluded, as were images flagged \texttt{iscrowd=1} or containing more than three person annotations. Retained bounding boxes were required to occupy between $5\%$ and $60\%$ of the frame area and to sit at least $2\%$ away from the image edge. They also needed a height-to-width ratio of at least $0.5$, to bias toward standing rather than incidental figures. Post-download quality checks additionally required a minimum resolution of $256\,\mathrm{px}$, a Laplacian blur variance of at least $100$, and a minimum saturation of $15$ to reject grayscale or infrared frames. This filtering and quality-check pipeline was implemented as a dedicated download script.  % scripts/download_coco_humans.py

Third, a dedicated Wikimedia Commons category-crawling pipeline was built for the species that remained persistently under-represented after the first two steps. It was designed as a human-in-the-loop instrument rather than as a downloader: the crawl produces a reviewable category tree per label, and only categories surviving review are enumerated and fetched. The automated filter that performs the bulk of that review encodes an explicit account of what a wildlife photograph is \emph{not}. That exclusion vocabulary is the substance of the design: artwork and illustration, anatomical specimens and isolated body parts, range maps, philatelic material, tracks and spoor, taxidermy and trophies, food products, fossils, and categories in which the animal is incidental to the subject (\enquote{things named after lions}, \enquote{people with cheetahs}). Exclusions cascade to subcategories, on the reasoning that Commons subcategories inherit their parent's framing: every subdivision of an art category is still art.

Partial-animal imagery (heads, paws, trophy shots) is attacked at three independent points: in this category vocabulary, in the caption-based rescue stage, and in the caption evaluation criteria of \Cref{sec:staged_filtering_pipeline}. None of the three is individually reliable, and the redundancy is deliberate.

\paragraph{A Divergence Between the Stated and Implemented License Policy}
An audit of the assembled corpus against the policy stated above found that the two do not agree. The license allow-list actually applied to the Wikimedia and COCO downloads admits ShareAlike variants, and in the COCO case NoDerivatives as well, rather than restricting to CC0, CC-BY and public domain as the policy specifies. Cross-referencing the retained images against their recorded license strings quantifies the consequence: of the $11{,}066$ Wikimedia images in the final splits, $6{,}222$ carry a ShareAlike license. This breaks down to $4{,}213$ in training, $223$ in validation, and $1{,}786$ in test. The models reported in \Cref{chapter:results} were therefore trained and evaluated on a corpus that contains ShareAlike-licensed imagery, contrary to the policy this section sets out.

This is recorded here rather than resolved, for two reasons. The remedy (removing the affected images and retraining every model) is a substantial undertaking that falls outside the scope of the present work. The residual risk is a legal question about whether trained weights are a derivative work, which is precisely the unsettled question that motivated the conservative policy in the first place. The finding is carried forward as a named limitation (\Cref{sec:limitations}) and, for any production use of these weights, as a prerequisite to be settled before deployment rather than a matter of academic record.

%%%%%%%%%%%%%%%%%%%%%%%%%%%%%%%%%%%%%%%%%%%%%%%%%
\subsection{Resulting Corpus Composition}\label{sec:corpus_composition}

The sourcing effort described above produced a raw, multi-source corpus spanning GBIF and iNaturalist observations, Open Images and COCO subsets, and Wikimedia Commons crawls, covering $216$ of the $225$ target classes with at least one image. \Cref{tab:source-funnel} gives its composition after the quality filtering of \Cref{sec:data_quality_and_curation} has been applied. The per-source survival rate is itself informative: it varies from $53\%$ to $85\%$, and the ordering matches the trust assignment made on licensing and provenance grounds above. The community-curated observation platform survives best, and the two unverified sources survive worst. Of the $465{,}242$ images passing quality filtering, $457{,}468$ are attributable to one of the $225$ target classes. The remainder failed taxonomic matching and are not usable.

\begin{table}[H]
\centering
\caption{Per-source composition of the raw corpus and its survival through quality filtering.}
\label{tab:source-funnel}
\begin{tabular}{@{}lrrrr@{}}
\toprule
Source & Passed & Failed & Total & Pass rate \\
\midrule
iNaturalist & $400{,}067$ & $73{,}443$ & $473{,}510$ & $84.5\%$ \\
GBIF & $39{,}388$ & $27{,}493$ & $66{,}881$ & $58.9\%$ \\
Wikimedia Commons & $12{,}598$ & $5{,}760$ & $18{,}358$ & $68.6\%$ \\
Open Images & $7{,}688$ & $4{,}059$ & $11{,}747$ & $65.4\%$ \\
Stock imagery & $5{,}501$ & $4{,}820$ & $10{,}321$ & $53.3\%$ \\
\midrule
Total & $465{,}242$ & $115{,}575$ & $580{,}817$ & $80.1\%$ \\
\bottomrule
\end{tabular}
\end{table}

Coverage at this point remained heavily uneven despite the volume: $83$ classes had at least $1{,}500$ images, while $29$ had fewer than $100$ and a further $58$ fewer than $500$. Two further sources contribute imagery that is not wildlife at all: the COCO-derived human class described above, and $174$ background photographs containing no animal. These background images are retained as true negatives, so that the detector sees frames in which the correct answer is to predict nothing. A third is not an image source at all: species articles were scraped from Wikipedia to serve as the textual corpus from which synthetic-image prompts are built (\Cref{sec:synthetic_generation_pipeline}).

This raw corpus is not yet fit for training: it mixes sources of very different reliability, contains mislabeled and low-quality images, and has not yet been split into train, validation, and test partitions. \Cref{sec:data_quality_and_curation} describes the staged quality-filtering and curation pipeline that turns this raw pool into the curated dataset used for training, and \Cref{sec:synthetic_data_supplementation} describes how the small number of classes that remained data-poor even after sourcing were supplemented with synthetic imagery.
